\section{Finite real subdivision and passage words}
\label{sec:subdivision}

We work with all coefficient vectors of polynomial fields $V=(P,Q)$ of
degree at most $d$.  The spatial domain is a fixed bounded square.
There is no compactness assumption on the coefficient vector.  A dilation
puts any prescribed finite collection of periodic orbits in this square,
without increasing the degree.  All constants and all finite lists below
are allowed to depend on $d$ and on previously fixed subdivision choices.
They are independent of the field and of its periodic orbits.

The finite lists below consist of formulas and labels whose parameter
values vary with the field.  An orbit selects a word in this alphabet
and its endpoint values.  The further factorizations used for an
absolute-zero test are chosen only after a sequence is fixed and do
not enlarge the alphabet.

\subsection{Preparation, clocks, and small arguments}

Here is the precise external preparation result we use.  Coordinates $p$
in this statement include all parameters; $x$ is the last variable.

\begin{theorem}[Simultaneous globally subanalytic preparation]
\label{thm:subdivision-preparation}
Let $X\subset\RR^k\times\RR$ and $f_1,\ldots,f_N:X\to\RR$ be
globally subanalytic.  There is a finite globally subanalytic cell
decomposition of $X\cap\{x\ne0\}$ such that on each cell the functions
have a common center $a(p)$ and representations
\[
 f_\nu(p,x)=c_\nu(p)|x-a(p)|^{q_\nu}
 U_\nu\bigl(b_{\nu j}(p)|x-a(p)|^{r_{\nu j}}:1\le j\le h_\nu\bigr).
\]
All exponents are rational, all displayed parameter functions are globally
subanalytic, and each $c_\nu$ is identically zero or nowhere zero.
The sign of $x-a(p)$ is constant.  Each argument tuple lies in
$[-1,1]^{h_\nu}$, and $U_\nu$ is real analytic on a neighborhood of that
cube and positive there.  Moreover, either $a=0$ or
$|x-a(p)|<\epsilon|x|$ on the cell for some $\epsilon<1$.
\end{theorem}

This is the globally subanalytic case of Lion--Rolin preparation:
\cite[\S0.3, Theorem~1 in \S0.4, and \S1.1(5)]{LionRolin1997};
the stated finite-family formulation is
\cite[Definition~2.14 and Fact~2.15]{Opris2023}.
The graph $x=0$ is separated before using the theorem.  The last
center inequality will not be needed.  A finite union of argument lists
makes the arguments common as well.  Exponent-zero arguments are bounded
parameter arguments.  We freely split signs and zero loci.

We also use finite analytic cell decomposition and uniform finiteness for
globally subanalytic families.  The analytic refinement can be obtained
within the cited preparation theorem: recursively refine the base for
all its finitely many coefficients, centers, and cell-boundary functions.
On every fixed-sign interval the rational powers and strong units are
analytic.  Induction on the number of coordinates, putting parameters
first, therefore gives compatible analytic cylindrical pieces; see also
\cite[\S1.1(1),(6)]{LionRolin1997}.  Uniform bounds for fiber components
are supplied by \cite[Theorem~3.14]{BierstoneMilman1988}, after algebraic
compactification for globally subanalytic sets.  In particular, a fixed such family of
sets of fiber dimension at most one admits a uniform finite decomposition
of its fibers into points and injectively parametrized regular analytic
arcs.  Adding finitely many globally subanalytic functions allows their
regularity and signs to be fixed on these pieces.  Differentiation on a
regular piece, algebraic root selection, finite maxima, and positive
rational powers preserve global subanalyticity.  These facts concern the
preparation data on their regular pieces; they assert no uniform bound on
their derivatives.

Every strong unit has uniform real upper and positive lower bounds.
Compactness of its argument cube also gives a complex neighborhood on
which it is holomorphic and nonvanishing, after decreasing that
neighborhood.  If a prepared function is bounded, its leading monomial
is bounded, since its unit is bounded below.  We may therefore include
the leading monomial as an additional bounded argument whenever only a
bounded analytic representation is needed.  Real and imaginary parts of
complex algebraic data are prepared separately.

\begin{lemma}[Compatible change of Euler center]
\label{lem:subdivision-clock}
Suppose that the old clock is $D_0=(x-a_0(p))\partial_x$.
After simultaneous preparation and finite subdivision, all the required
data can be represented in an Euler clock
\[
 D=(x-a(p))\partial_x=\tau\partial_\tau=hD_0,
 \qquad \tau=|x-a(p)|>0,\qquad 0<|h|\le2.
\]
If the initial clock is $\partial_x$, the same assertion holds without
the multiplier bound at this first step.  Previously used time data can
be included in the preparation.
\end{lemma}

\begin{proof}
For a center $a$ furnished by preparation, keep $a$ on
$|x-a|\le2|x-a_0|$.  The multiplier is
$h=(x-a)/(x-a_0)$.  Put all zero denominators and centers on the
boundary.  On the complementary region return to $a_0$.  With
$q=(x-a_0)/(a_0-a)$ the complementary inequality is
$|1+q|>2|q|$, which implies
\[
 -\tfrac13<q<1,\qquad \tfrac23<1+q<2.
\]
Consequently, for every rational $r$ occurring in the finite list,
\[
 |x-a|^r=|a_0-a|^r(1+q)^r.
\]
The second factor is analytic with a fixed complex neighborhood of the
indicated real interval.  If $b|x-a|^r$ was bounded, then
$b|a_0-a|^r$ is bounded, with a uniform bound depending only on $r$.
The resulting analytic compositions have compact argument ranges with
room.  Now $h=1$.  Splitting its sign, when necessary, finishes the proof.
\end{proof}

The next observation supplies the small pure-time profiles used in
Section~\ref{sec:terminal}.  Suppose finitely many bounded arguments
$b_j\tau^{r_j}$ occur, with nonzero $r_j$.  Fix a desired smallness
$\delta>0$.  On pieces where all are smaller than $\delta$, keep the
clock.  Otherwise select one with
$\delta\le|b_j\tau^{r_j}|\le M$.  With
$t=|b_j|^{-1/r_j}$, the ratio $\tau/t$ belongs to a fixed compact
subinterval of $(0,\infty)$.  Subdivide that interval into finitely many
relative bins.  Choose a baseline $\tau_*=ct$ slightly below each bin,
so that throughout it
\[
 \tau=\tau_*+\widetilde\tau,\qquad
 0<\widetilde\tau/\tau_*<\eta.
\]
All old powers become constant multiples of
$(1+\widetilde\tau/\tau_*)^r$.
Previously bounded multipliers $b_j\tau_*^{r_j}$ remain bounded.
A small mixed argument $b\tau^a u^s$, $s\ne0$, becomes
$b\tau_*^a u^s(1+\widetilde\tau/\tau_*)^a$ and remains small after
an arbitrarily small fixed adjustment to its tolerance.  The new clock
multiplier is $\widetilde\tau/\tau$, and
\begin{equation}
\label{eq:subdivision-contact-power}
 \frac{\widetilde\tau}{\tau}A\tau^\lambda
 =A\tau_*^{\lambda-1}\widetilde\tau
 (1+\widetilde\tau/\tau_*)^{\lambda-1}.
\end{equation}
Thus these \emph{contact bins} have leading time exponent $1$.
Choose finitely many analytic neighborhoods of the compact tuples of
baseline constants first, and then make $\eta$ small inside their
common margins.  This ensures that setting the new small arguments to
zero stays within the analytic neighborhoods.  Narrowing bins changes
their number, but introduces no new exponent other than the patterns
in Equation~\eqref{eq:subdivision-contact-power}.

\subsection{Poles, annuli, and pole-free boxes}

Away from $P=0$ the scalar equation is $\partial_x z=Q(x,z)/P(x,z)$.
We need a slightly more general construction, since it will also be
used after rational weighting: start with
\begin{equation}
\label{eq:subdivision-rational}
 D_0z=R(p,x,z),
\end{equation}
where $R$ is rational in $z$ of bounded degree and its base coefficients
are globally subanalytic and analytic on the selected base pieces.
The clock is either the first clock $\partial_x$ or an Euler clock.

Partition by denominator degree, root multiplicities, and coefficient
regularity.  The complex roots $p_j(p,x)$ can be enumerated on finitely
many regular pieces, with multiplicities.  Choose a pole whose real part
$c$ is horizontally nearest to the real state $z$, and split by this
choice and by the sign of $z-c$.  Set $u=\sigma(z-c)>0$,
$\sigma\in\{1,-1\}$.  For every denominator root,
\begin{equation}
\label{eq:subdivision-pole-distance}
 |z-p_j|\ge |z-\Re p_j|\ge |z-c|=u.
\end{equation}
With no denominator roots use $c=0$ instead.  The transformed equation is
$D_0u=\sigma R(p,x,c+\sigma u)-\sigma D_0c$.
Its numerator remains polynomial in $u$ after using the denominator;
the degree is uniformly bounded.  In particular the derivative of the
center is retained exactly.  A numerator that vanishes identically is a
constant-state passage.  Degree drops and zero coefficients are separated
before any ratios are formed.  Common numerator and denominator factors
cause no problem: retaining a removable pole only adds a boundary.

Write the translated denominator factors as $u-\varpi_j$, where
$\varpi_j=\sigma(p_j-c)$, and the numerator as $\sum_{l=0}^{N}n_lu^l$.
The \emph{complete radius list} for the following split is
\begin{equation}
\label{eq:subdivision-complete-radii}
 \mathcal R=
 \{\rho_j=|\varpi_j|:\varpi_j\ne0\}
 \ \cup\ 
 \{|n_l/n_r|^{1/(r-l)}:0\le l<r\le N,\ n_ln_r\ne0\}.
\end{equation}
In particular every positive translated pole modulus is a radius,
whether or not any two numerator terms are present.  Zero displacements
are separated and will contribute exact powers of $u$.
Prepare simultaneously this list, the nonzero coefficients, the real
and imaginary parts of the bounded directions $\varpi_j/\rho_j$, the
centers and their needed derivatives, and all earlier time-coordinate
data needed to return to physical endpoints.  Use
Lemma~\ref{lem:subdivision-clock}.  For every $\rho\in\mathcal R$
the preparation gives a positive pure representative
$s_\rho=C_\rho(p)\tau^{\gamma_\rho}$ and, for a single fixed $K\ge1$
on the finite family,
\begin{equation}
\label{eq:subdivision-radius-comparison}
 K^{-1}s_\rho\le\rho\le Ks_\rho.
\end{equation}
There are at most $q+N(N+1)/2$ entries, where $q$ is the denominator
degree.  Algebraic root selection, modulus, and positive rational
powers make them globally subanalytic on their regular pieces.  Repeated
or equal radii may be retained.  Thus including the pole moduli changes
neither the finite-family hypothesis of preparation nor degree control.

Choose a large fixed factor $H>K$.  The far domain is the union of the
pieces on which, for \emph{every} $\rho\in\mathcal R$, either
$u>Hs_\rho$ or $u<s_\rho/H$.  Fix these alternatives on each piece.
Equation~\eqref{eq:subdivision-radius-comparison} implies respectively
$\rho/u<K/H$ or $u/\rho<K/H$.  One numerator term therefore
dominates.  Indeed the ratio
of terms of degrees $l<r$ has absolute value
$(u/|n_l/n_r|^{1/(r-l)})^{r-l}$ or its reciprocal; the selected far
inequalities make all ratios to the largest term at most $K/H$ once
$H>K$.  For a nonzero pole displacement the exact extractions are
\[
 u-\varpi_j=u(1-\varpi_j/u)
 \quad\hbox{or}\quad
 u-\varpi_j=-\varpi_j(1-u/\varpi_j),
\]
and the chosen correction has modulus less than $K/H$.
The normalized directions have modulus one, so the bounds hold also
for nonreal poles and for several poles with equal real part.  Poles
at displacement zero contribute an exact power of $u$; conjugate
factors are paired to retain real formulas.
Thus these pieces have the exact form
\begin{equation}
\label{eq:subdivision-annulus}
 Du=A\tau^\lambda u^m F(X),\qquad
 X_j=b_j\tau^{a_j}u^{s_j},\qquad u>0,
\end{equation}
where $A\ne0$ is parameter-only, $m\in\ZZ$, and all other exponents
are rational.  Constant arguments have been included in a bounded
ordinary tuple, suppressed in this notation.

Here and below ``exact'' refers to equality of the scalar fields on the
piece, not to an asymptotic approximation.  Prepared units occur in the
small ratios just described and are retained in $F$.  After the contact
construction and division by the nonzero value at zero profiles, we have
$F(0)=1$ on larger analytic boxes and, for any prescribed fixed
$\varepsilon>0$, can arrange
\begin{equation}
\label{eq:subdivision-annulus-small}
 \max_j|X_j|<\varepsilon,\qquad
 |F-1|+|\partial_{\log u}F|<\varepsilon.
\end{equation}
For the second inequality use
$\partial_{\log u}F=\sum s_jX_jF_{X_j}$ and Cauchy bounds on the
larger boxes.  The leading-term and pole-extraction alternatives form a
finite list before $H$ is increased.  Their bounded unit arguments have
fixed analytic neighborhoods.  Hence the required smallness can be
chosen from this list before choosing the final far factor.  Contact
bins contribute bounded baseline constants and fixed power patterns, so
their later narrowing does not change that assertion.  In particular,
one may anticipate any fixed compact power-clock state boxes required
in Section~\ref{sec:terminal}, whose state factors then cost only fixed
constants.

On each remaining piece there is a radius $\rho\in\mathcal R$ with
$s_\rho/H\le u\le Hs_\rho$.  Choose one such radius by a finite
subdivision, set $s=s_\rho=C\tau^\gamma$, and put $z_1=u/s$.
Thus $H^{-1}\le z_1\le H$, whether the selected radius is a pole
modulus or a numerator balance.  Equalities in these cuts are boundary
pieces.  The equation is
\[
 Dz_1=s^{-1}Du-\gamma z_1.
\]
By Equation~\eqref{eq:subdivision-pole-distance} every pole in the
$z_1$-plane has distance at least $z_1\ge1/H$ from the actual real
point.  Choose fixed real bins of radius $\delta_1<1/(8H)$ and
concentric larger complex disks of radius $2\delta_1$.  For each bin
restrict the base to parameters for which the actual piece meets the
bin.  The triangle inequality from any such actual point gives pole
distance at least $1/H-3\delta_1>1/(2H)$ throughout the larger disk.
Thus the lower bound at actual points supplies the required complex
disk bound; it is not presumed to hold at arbitrary bin centers.
These finitely many projected base sets are globally subanalytic and
can be regularly subdivided.  On the disks the field is rational with
numerator degree at most $\max(N,q+1)$ and denominator degree $q$:
the new term $-\gamma z_1$ adds at most $z_1$ times the denominator.
This remains true when several pole radii are comparable, when some
other pole radii tend to zero or infinity relative to $s$, and when
several poles have the same nearest real part.

\begin{example}[A conjugate pole pair on the comparable branch]
\label{ex:subdivision-comparable-poles}
For $R=1/(z^2+\epsilon^2)$, $0<\epsilon<1/10$, both poles have real
part zero.  On $z>0$ we have $c=0$, $u=z$, and the radius list is
$\{\epsilon,\epsilon\}$, with no numerator balance radius.
The point $u=\epsilon$ belongs to the comparable branch $s=\epsilon$.
In the original clock,
\[
 z_1=u/\epsilon,\qquad
 \partial_xz_1=\frac{\epsilon^{-3}}{1+z_1^2}.
\]
Near $z_1=1$ the poles are $\pm\mathrm i$, so, for example, the disk
$|z_1-1|<1/2$ is uniformly pole-free.  On the smaller real interval
$|z_1-1|<1/4$, the normalized field is at least $16/41$ and bounded
on the larger disk.  This is a nozero box with scale comparable to
$\epsilon^{-3}$.  Changing to an Euler clock only adds its prepared
time factor.  A purported far-annulus description at $u=\epsilon$
would instead have $|u/\epsilon|=|\epsilon/u|=1$; the complete radius
list is precisely what sends these points to the box construction.
\end{example}

\begin{lemma}[Uniform box data]
\label{lem:subdivision-box-data}
After the preceding finite subdivision, each nonzero comparable-scale
field has, on a fixed larger state disk $\mathcal D$ and a smaller real
interval $I\Subset\mathcal D$, a positive globally subanalytic scale
$L(p,x)$ and constants $K,m$ such that
\begin{equation}
\label{eq:subdivision-jet}
 \sup_{z\in\mathcal D}|F(p,x,z)|\le KL(p,x),\qquad
 \sum_{j=0}^{m}|\partial_z^jF(p,x,z)|\ge L(p,x)
 \quad(z\in I).
\end{equation}
The normalized rational coefficients belong to bounded sets with uniform
analytic room.  The constants and degrees are fixed on each of finitely
many boxes.
\end{lemma}

\begin{proof}
Normalize the denominator by its value at the disk center.  In its
factored representation its coefficients and its reciprocal are bounded
on a slightly smaller disk, because all poles are uniformly separated.
Let $L_0$ be the maximum modulus of the coefficients of the resulting
numerator.  Separate $L_0=0$; otherwise divide the numerator by $L_0$.
The normalized numerator and denominator belong to compact coefficient
sets, the denominator is uniformly nonzero, and at least one numerator
coefficient has modulus one.  If the numerator degree is at most $m$,
no member can have its first $m+1$ jets vanish at a point: multiplication
by the nonzero denominator would give a polynomial of degree at most
$m$ with a zero of order $m+1$.  Compactness, including $z\in\overline I$,
gives a positive uniform lower bound on the sum of those jets.  Rescaling
$L_0$ by that fixed bound gives Equation~\eqref{eq:subdivision-jet}.
The supremum estimate follows from the coefficient bounds.  All the
normalizations use subanalytic operations and nonzero divisors.
\end{proof}

\subsection{The finite box induction}

The order $m$ in Equation~\eqref{eq:subdivision-jet} is the induction
index.  During this argument the state disks can be decreased by fixed
amounts and then covered by finitely many fixed smaller disks.  A
``uniform bound'' always refers to such specified larger disks.

A box may also carry an absolute bound for its field on its larger state
disk.  Every nonterminal child below either has lower jet order or
carries such a bound; once the bound is present, the next continuing step
lowers the order.  At each selected index, the possible center branches and their
preparations will be fixed before the final lower-jet thresholds.

Choose positive thresholds
$\epsilon_0,\ldots,\epsilon_m$ with
$\sum\epsilon_j<1$.  Their choice is made from higher to lower indices,
and lower ones will be made sufficiently small below.  Assign a point
to the least $i$ with $|\partial_z^iF|\ge\epsilon_iL$; some such $i$
exists by Equation~\eqref{eq:subdivision-jet}.  Split its sign and the
signs of the lower derivative expressions on regular cells.  Cauchy
estimates on the larger disk bound every fixed further derivative by
a constant times $L$.

If $i=0$, the field has fixed nonzero sign and a lower bound on the piece.
If $i=1$, the smallness of $|F|/L$ relative to the derivative margin
places the point as close as desired to a unique real simple root.
For completeness, choose a radius $r_0$ so that the derivative changes by
at most half its lower bound on the corresponding real interval and
complex disk.  If $|F(z)|<r_0\epsilon_1L/4$, the values at
$z\pm2|F(z)|/(\epsilon_1L)$ bracket zero.  Strict monotonicity gives a
unique real root $c$, and $|F_z(c)|\ge\epsilon_1L/2$.
The analytic implicit theorem gives a root chart with a fixed smaller
complex neighborhood.  This proves both existence and the needed room.

Prepare $L$ and the bounded normalized rational coefficients in these
two cases, together with the coordinate data, and perform the final
contact refinement.  The resulting equation is
\begin{equation}
\label{eq:subdivision-terminal-box}
 Dz=B\tau^\lambda f(X,z),\qquad
 X_j=b_j\tau^{r_j},\qquad B>0,
\end{equation}
where $r_j\ne0$ are rational, $X$ is as small as required, and $f$ is
analytic with uniform bounds on larger boxes.  Constant arguments are
bounded ordinary parameters.  Either $f$ has constant sign with a fixed
positive lower bound on the piece, or the state is arbitrarily close to
a regular root $z=c(X)$, with a fixed derivative margin.  In the latter
case the root at $X=0$ is still regular.  To see that these choices do
not depend circularly on the final root-tracking tolerance, first prepare
the bounded coefficient data on the full pole-free bins.  The derivative
margin fixes the root neighborhoods there.  Only then choose the profile
smallness and the order-zero threshold.  Finitely many ordinary root
charts cover the compact normalized coefficient sets.

Now let $i\ge2$.  Apply the same argument to $\partial_z^{i-1}F$.
There is a nearby real root $c(p,x)$ with
\[
 \partial_z^{i-1}F(p,x,c)=0,\qquad
 |\partial_z^iF(p,x,c)|\ge\tfrac12\epsilon_iL.
\]
Regular root selection and subdivision make $c$ analytic in $x$ and
globally subanalytic.  Its distance from the actual state can be made
arbitrarily small by the lower thresholds.  Set $r=z-c$ and write the
exact translated field as
\begin{equation}
\label{eq:subdivision-center}
\begin{gathered}
 D_0r=A_0+\sum_{l=1}^{i}A_lr^l+A_ir^{i+1}H(p,x,r),
 \\
 A_0=F(p,x,c)-D_0c,\qquad
 A_l=\frac{\partial_z^lF(p,x,c)}{l!}\ (l>0).
\end{gathered}
\end{equation}
Here $A_{i-1}=0$, $A_i\ne0$, and $H$ has a uniform analytic bound
on a fixed disk.  It is rational analytic there, with bounded normalized
subanalytic coefficient data: subtract the finite Taylor polynomial in
the normalized rational representation.  For $1\le l<i-1$,
$|A_l/A_i|$ can be made arbitrarily small.  Indeed the lower derivatives
are small at the original point, the center displacement is small, and
Taylor's formula with the above Cauchy bounds transfers these estimates
to $c$.  No bound has been asserted for $D_0c$ or for $A_0$.

Define
\begin{equation}
\label{eq:subdivision-radius}
 S=\max_{0\le l<i}|A_l/A_i|^{1/(i-l)}.
\end{equation}

To make the exponent list independent of the final lower-jet thresholds,
prepare the full center branches before imposing those thresholds.
Fix the current disks, normalized rational coefficient bounds,
and $\epsilon_i$.  Choose a compact real interval $J$ containing the
actual smaller state interval with room and lying a fixed distance
inside the larger disk.  The \emph{full candidate-center set} is
\begin{equation}
\label{eq:subdivision-candidate-centers}
 \mathcal C_i=
 \left\{(p,x,c):c\in J,\quad
 \partial_z^{i-1}F(p,x,c)=0,\quad
 |\partial_z^iF(p,x,c)|\ge\tfrac12\epsilon_iL(p,x)\right\}.
\end{equation}
No condition involving a lower-jet threshold is imposed here.  Since
$F$ is rational of fixed degree, clearing the nonvanishing denominator
gives a polynomial equation in $c$ of fixed degree.  Its roots in
$\mathcal C_i$ are simple by the displayed derivative margin.
The identically zero polynomial has no roots satisfying that margin.
Thus finite subanalytic regular subdivision supplies finitely many
candidate root branches.  Take the compact ambient normalized
coefficient set of Lemma~\ref{lem:subdivision-box-data}, including the
closure of the attainable coefficient image; its denominator remains
uniformly nonzero on the state disk.  In its product with $J$, impose
the root equation and the closed derivative margin for the normalized
field $F/L$.  This is a compact set, even when the attainable parameter
image itself is not closed.  The analytic implicit theorem on this
compact ambient root set provides the root charts and their uniform
neighborhoods before any final lower-jet restriction.

On every full branch form all $A_l$, including
$A_0=F(c)-D_0c$, and $S$ from
Equation~\eqref{eq:subdivision-radius}.  Their base functions are
globally subanalytic; differentiating the root equation on its regular
base gives $c_x=-\partial_x\partial_z^{i-1}F/\partial_z^iF$, so this
claim neither assumes nor supplies a bound on $c_x$.
Prepare these full-branch functions, their bounded normalized analytic
data and the earlier coordinate data, using
Lemma~\ref{lem:subdivision-clock}.  Every $A_l$ is multiplied by the
same $h$, so $S$ is unchanged.  On each prepared piece with $S>0$ write
$s=C\tau^\gamma\asymp S$ and fix $K_0\ge1$ such that
$K_0^{-1}S\le s\le K_0S$.  The possible exponents $\gamma$ and
comparison constants are now fixed before the lower-jet restrictions.

Fix a small cutoff $\delta$ for $S$.  Arrange the positive lower
coefficient radii to be less than $\delta/2$.  On $S\ge\delta$ the
maximum in Equation~\eqref{eq:subdivision-radius} must come from $A_0$.
On $|r|\le\eta\delta$, with fixed sufficiently small $\eta$, every
positive-degree term in Equation~\eqref{eq:subdivision-center} is a
small fraction of $|A_0|$.  The remainder has the same property.  The
actual state lies in a still smaller such disk by the choice of lower
thresholds.  This is a nozero box with scale $|hA_0|$.

It remains to consider $S<\delta$.  If $S=0$, discard $r=0$ as a
boundary and use the monomial field of power $i$.  If $S>0$ and
$|r|>Hs$, with $H$ fixed sufficiently large, then
\[
 |(A_l/A_i)r^{l-i}|\le (S/|r|)^{i-l}
 \le (K_0/H)^{i-l}.
\]
The higher remainder is small by the smallness of $|r|$.  Thus this is
again an annulus, now of state power $m=i\ge2$, with the prepared unit
and contact refinements already described.

For the inner part put $r=sz_2$, $|z_2|\le H$.  On a fixed larger disk
we have exactly
\begin{equation}
\label{eq:subdivision-inner}
\begin{gathered}
 Dz_2=L_s\left(z_2^i+\sum_{l<i-1}d_lz_2^l+E(p,x,z_2)\right)
 -\gamma z_2,\\
 L_s=hA_is^{i-1},\qquad
 d_l=(A_l/A_i)s^{l-i}.
\end{gathered}
\end{equation}
where $\|E\|\le C s$.  Its fixed finite jets on smaller disks have the
same bound.  All $d_l$ are bounded, and at least one is bounded away
from zero: choose an index attaining $S$ in
Equation~\eqref{eq:subdivision-radius}, and use $S/s\in[K_0^{-1},K_0]$.
In particular the polynomial coefficient set in
\[
 P_0(t)=t^i+\sum_{l<i-1}d_lt^l
\]
is compact and excludes $t^i$.

There are three cases, with fixed thresholds $0<a<b$.
\begin{enumerate}
\item If $|\gamma/L_s|<a$, including $\gamma=0$, the normalized jet
order drops to $i-1$.  To prove it, an $i$-fold zero of $P_0$ would give
$P_0(t)=(t-t_0)^i$.  The missing coefficient of $t^{i-1}$ forces
$t_0=0$, contrary to its nonzero lower coefficient.  Compactness in
both coefficients and $|t|\le H$ therefore gives a positive lower bound
for $\sum_{j<i}|P_0^{(j)}(t)|$.  Choose $a$ and then $\delta$ so small
that $-(\gamma/L_s)t$ and $E$ preserve half this lower bound.  The
supremum bound is uniform as well, with scale $|L_s|$.
\item If $|\gamma/L_s|>b$, divide the field by $|\gamma|$.  Its state
derivative is $-\sgn(\gamma)+O(1/b)$ on the smaller disk.  Increasing
$b$ gives a first-jet lower bound and a uniform supremum bound.  Its
order is therefore at most one.
\item If $a\le|\gamma/L_s|\le b$, retain order $i$.  The $i$th
derivative of the normalized polynomial is $i!$ and the remainder is
small.  Moreover $\gamma\ne0$ and $|L_s|\le|\gamma|/a$.  Since the
possible $\gamma$ form a finite fixed list, the \emph{unnormalized}
field in Equation~\eqref{eq:subdivision-inner} has an absolute uniform
bound on its larger disk.
\end{enumerate}

The third case can occur at most once before the order decreases.
To verify this, carry that absolute disk bound to the next step.  If
its selected index is less than $i$, the order has already decreased.
If the selected index is $i$, its Taylor coefficient $A_i$ at the next
center is absolutely bounded by Cauchy's inequality.  For each of the
finitely many new possible nonzero exponents $\gamma'$, choose the new
cutoff so small that
\[
 |h'A_i(s')^{i-1}|<|\gamma'|/b.
\]
More explicitly, if the inherited bound is $M$, the centers have disk
margin $\rho$, the new comparison is $s'\le K_0'S'$, and
$\gamma_*'$ is the minimum of the finitely many nonzero
$|\gamma'|$, it suffices to impose
\[
 2M\rho^{-i}(K_0'\delta')^{i-1}<\gamma_*'/b.
\]
There is no condition of this kind if the nonzero exponent list is
empty.  This is possible since $i-1>0$.  The inner child is
then in the second case.  If $\gamma'=0$, it is in the first case.
The large-$S'$ and outer children are already terminal nozero boxes or
annuli.  Thus a second retention at order $i$ is impossible.

The choices can now be made in the required order.  After the full-branch
preparations, choose the outer factor $H$ and then the thresholds $a,b$
from the resulting compact coefficient bounds.  Next choose the cutoff
$\delta$ required for the compact inner disks, small remainders and,
when present, the inherited absolute bound.  Finally choose the lower
jet thresholds small enough to give the center displacement and the
large-$S$ nozero estimate.  Any subsequent restriction only restricts
these prepared identities and introduces no new exponent.  Children of
smaller order make their own finite choices.

We have proved the following assertion, with a genuine finite recursion:
at a selected index $i$, every child is terminal, has smaller order,
or has order $i$ together with an absolute bound which forces its next
nonterminal child to have smaller order.  The order starts finite and
is a nonnegative integer.  Every node has finitely many children by
preparation and cell decomposition.  Consequently the whole tree is
finite.  Equivalently, label a box by $(m,e)$, with $e=0$ when an
inherited absolute bound is available and $e=1$ otherwise.  A retention
changes $(m,1)$ to $(m,0)$; every other continuing child has smaller
order.  The nonnegative integer $2m+e$ strictly decreases.

\subsection{One weighted redo at a nonzero resonance}

Only annuli from the initial pole split can have state power $m=1$.
Annuli created in the box induction have power at least two; contact
bins have time exponent one.  The remaining special case for
Section~\ref{sec:terminal} is therefore
\[
 Du=A_0uF(X),\qquad X_j=b_j\tau^{a_j}u^{\beta_j}.
\]
All pairs $(a_j,\beta_j)$ are rational and nonconstant.  The nonzero
generator resonances are the finite set
$\alpha=-a_j/\beta_j\ne0$, for $\beta_j\ne0$.

\begin{lemma}[Termination of the resonant redo]
\label{lem:subdivision-resonance-redo}
Around these finitely many resonances, apply the rational weighting
$U=u/\tau^\alpha$ and repeat the preceding rational subdivision once.
The old resonance widths and old generator smallness can be chosen so
that every new annulus with state power one and time exponent zero has
a slope in its own small-slope regime.  In particular no second nonzero
resonance redo is required.
\end{lemma}

\begin{proof}
In the old Euler clock the exact weighted equation is
\[
 DU=E(p,x,U)=U(A_0F-\alpha),\qquad
 E_U=A_0F-\alpha+A_0\partial_{\log u}F.
\]
On a sufficiently narrow fixed neighborhood of $A_0=\alpha$ and with
Equation~\eqref{eq:subdivision-annulus-small} sufficiently strong,
$|E_U|<\varepsilon$, for any predetermined $\varepsilon>0$.
The field $E$ is rational in $U$, with globally subanalytic base
coefficients: rational weighting changes only its base coefficients.

First perform the initial preparations of all these weighted rational
fields on their unrestricted domains.  They depend on the old exact
rational field, its uncontacted clock and state shift, and the finite
candidate weights.  They do not depend on any width around a resonance.
This gives a finite list of new initial uncontacted exponent pairs.
For every such list fix a positive small-slope threshold $c_*$.  One
possible choice, decreased by a fixed factor to allow units near one,
is
\[
 c_*<\min_{a_j\ne0}\frac{|a_j|}{8(1+\max_j|\beta_j|)};
\]
when the minimum is empty, take any fixed small positive threshold.
Then $|c|<c_*$ gives a uniform nonzero sign margin for all slopes
$a_j+\beta_jc$ with $a_j\ne0$ and excludes all nonzero resonances.
Take the minimum of these finitely many thresholds, and choose
$4\varepsilon$ below it.  Choose the original resonance and generator
thresholds to ensure $|E_U|<\varepsilon$.

Consider a new initial unscaled annulus.  The state change before
extracting its monomial is only $r=\sigma(U-c(p,x))$.  For a new Euler
clock $D'=hD$, Lemma~\ref{lem:subdivision-clock} gives $|h|\le2$, and
\[
 D'r=h\sigma(E-Dc),\qquad \partial_r(D'r)=hE_U.
\]
The unbounded derivative of the moving center cancels from the state
derivative.  If this annulus has $m'=1$, $\lambda'=0$, its field is
$A'_0rF'(X')$.  Its own subsequent far and generator cuts arrange
\[
 |F'+\partial_{\log r}F'-1|<\tfrac12.
\]
It follows that $|A'_0|\le4\varepsilon<c_*$.  The estimate is needed
only at the actual points of the restricted piece; the prepared
identities themselves were obtained on the unrestricted family.
All other children are boxes, box-induction annuli of power at least
two, or contact bins of time exponent one.  They cannot initiate a
nonzero linear resonance redo.  The new analytic smallness choices
come after the finite exponent thresholds and preserve their sign
margins.  This proves the asserted termination without a circular
choice of resonance widths.
\end{proof}

\subsection{Boundary crossings and uniformly finite words}

All cuts constructed so far are globally subanalytic in the physical
coordinates and polynomial coefficient parameters.  The same is true
of the additional real magnitude, state, root, and resonance cuts of
Section~\ref{sec:terminal}: they use rational powers on sign cells,
bounded analytic operations, and fixed comparisons.  The one-redo
lemma makes their recursion finite.  Endpoint choices between
alternative passage rules need no planar cut by an unbounded logarithm.

Fix the parameters.  By a boundary we mean the actual frontier in the
physical plane of the relevant open identity pieces.  Equalities which
hold on a whole open parameter fiber are not thin boundaries in that
fiber.  Globally subanalytic cell decomposition, applied with the
parameters first, separates these possibilities and ensures that the
actual boundary fibers have dimension at most one.  Include $P=0$,
all exceptional vertical base locations, and all zero state coordinates
excluded during changes of variables.  Split the boundary fibers into
uniformly finitely many points and simple regular analytic arcs.
Further split an arc, with regular oriented tangent $T$, by the signs
of $\det(V,T)$ and by the zeros of $V$.  This is another finite uniform
subanalytic subdivision.

The transverse-arc assertion below is the no-contact case of the planar
Rolle principle for separating trajectories
\citep[Lemma~1]{Khovanskii1984Rolle}. We include its Jordan-curve proof;
the tangential-arc assertion follows from uniqueness of the flow.

\begin{lemma}[Crossings of a boundary arc]
\label{lem:subdivision-crossing}
A periodic orbit meets each regular boundary arc on which
$\det(V,T)$ has a fixed nonzero sign at most once.  A connected
nonsingular boundary arc on which this determinant vanishes identically
is contained in one trajectory and hence in at most one periodic orbit.
\end{lemma}

\begin{proof}
A nonconstant periodic orbit of a locally unique planar flow, taken
with its least period, is an embedded Jordan curve: a repeated point
before that period contradicts uniqueness and minimality.  Its
orientation is the orientation of $V$.  Traversing a fixed oriented
transverse arc, consecutive crossings of a Jordan curve alternate
between entry into and exit from its bounded complementary component.
Their oriented crossing signs therefore alternate.  Here the crossing
sign is the fixed sign of $\det(V,T)$ (up to one fixed convention), so
two consecutive crossings are impossible.  Crossings are isolated by
transversality; on any compact subarc they are finite.  If there were
two crossings on the full arc, a compact subarc containing them would
already give the contradiction.

On an identically tangential nonsingular arc, $V$ is a continuous
nonzero scalar multiple of its regular tangent.  Reparametrizing the
arc therefore gives a solution of the field locally.  Uniqueness makes
the whole connected arc part of a single maximal trajectory.  Distinct
periodic trajectories cannot share it.
\end{proof}

Arcs on which $V=0$ cannot meet a periodic orbit.  The finite point
pieces include isolated tangencies, endpoints and singular points of
the boundary.  An embedded periodic orbit visits each such point at
most once.  By Lemma~\ref{lem:subdivision-crossing}, the number of
periodic orbits containing any nonsingular tangential arc is at most
the uniform number of these arcs.  Remove those finitely many orbits
from consideration.  Every remaining orbit has a uniform bound on its
number of boundary events, obtained by adding the numbers of point and
transverse arc pieces.  This argument concerns periodic Jordan curves;
no bound on the number of crossings of a nonclosed trajectory is used.

Near each event the field is nonsingular.  One of its physical
components is nonzero, so a sufficiently short piece of the orbit is a
regular graph in that physical clock.  Choose these neighborhoods
disjoint and their endpoints in the adjacent interior pieces.  This
produces one short graph connector per event.  The complementary
closed subarcs are compactly contained in their chosen identity pieces;
their positive transformed clock endpoints are finite and strictly
positive.  The connector neighborhoods can be as small as needed for
the particular orbit.  Their number is already uniformly bounded.

A nonconstant periodic orbit cannot avoid every graph-clock boundary:
its physical coordinate $x$ attains a maximum and a minimum, so $P=0$
somewhere on it.  If $x$ is constant on the entire orbit, its image lies
on a line and cannot be a Jordan curve.  This also disposes of fibers
with no useful nonvertical clock.  An orbit contained in the boundary
has a subarc in one of the finitely many regular pieces and is already
among the tangential exceptions.

\begin{theorem}[Finite passage templates]
\label{thm:finite-templates}
For each degree bound $d$, there is a finite collection of scalar
identity pieces, coordinate charts and passage rules, and integers
$N_d,E_d$, with the following properties.
\begin{enumerate}
\item On the square, every polynomial field of degree at most $d$ has
the stated finite globally subanalytic subdivision.  Its interior
scalar fields are constant-state fields, nozero or simple-root boxes
of Equation~\eqref{eq:subdivision-terminal-box}, or monomial annuli
of Equation~\eqref{eq:subdivision-annulus}.  Their analytic factors
have uniform larger neighborhoods and their nonconstant arguments
have the prescribed fixed smallness.  All changes back to physical
coordinates use parameter-only data, positive rational powers, and
bounded analytic operations.
\item The further real passage refinements in
Section~\ref{sec:terminal} are finite; in particular
Lemma~\ref{lem:subdivision-resonance-redo} allows at most one weighted
redo at a nonzero generator resonance.
\item Apart from at most $E_d$ periodic orbits containing tangential
boundary arcs, every periodic orbit is a cyclic concatenation of at
most $N_d$ interior passages and physical graph connectors.  There are
only finitely many resulting words, including all choices of chart,
sign, orientation, and the finitely many endpoint passage rules.
\end{enumerate}
The choices and bounds include all coefficient degeneracies.
\end{theorem}

\begin{proof}
The initial rational split is finite, and the box induction has finite
depth as proved above.  Zero numerators, vanishing coefficients and
degree changes were separated before normalization; multiple poles
were allowed throughout.  Thus the first assertion follows without
discarding a special coefficient fiber.  The real refinements of the
terminal cases use finite sign and magnitude alternatives, and their
only recursive operation is the redo covered by
Lemma~\ref{lem:subdivision-resonance-redo}.  The preceding boundary
argument gives $E_d$ and bounds the number of connectors and interior
links.  Each link has a label in a fixed finite alphabet, so words of
length at most that bound form a finite set.

Any subdivision into a number of normalized time fractions chosen
later for one sequence in the absolute-zero argument is an analytic
factorization of a given kernel, not an extra geometric link.  It does
not enter $N_d$ or the alphabet.  Finally, the parameter-only data here
are fixed for a field and a word; endpoint-dependent amplitudes and
normalizations belong to the endpoint graph variables of
Section~\ref{sec:counting}.  Thus no choice in this construction depends
on an unknown orbit cut when its value is declared parameter-only.
\end{proof}
